\section{Case Study}\label{sec:case-study}

The propositions of Section~\ref{sec:logic} are stated for a general GALA code, where the group data is the only input.  This section works one instance out end to end, so that every object in that section (the hypercube organization, the dirty cyclic shifts, and the fold-transversal Cliffords) appears as an explicit matrix.  We deliberately pick a small code, so that every logical action can be written out in full.  The instance is moreover \emph{strictly} self-dual, with the identity permutation as its ZX-duality, so its fold-transversal Clifford gates require no atom rearrangement at all, whereas a generic ZX-dual GALA code must first implement the fold of Prop.~\ref{def:dual-gala}.  Section~\ref{ssec:case-fold} makes this precise.

Let the code be $\mathcal{Q}\in\mathrm{GALA}_{12,5}(e\times C_{11})$ with $L=12$, $J=5$, where the top is trivial, of degree $1$, and the bottom $C_m=C_{11}$ is cyclic, so that $n=12\times1\times11=132$.  Every lift is a single group element, so the monomial hypotheses of Prop.~\ref{prop:distance-bound} apply: $J=5>L/4=3$, and the $J>L/4$ branch caps the distance at $d\leq L$.  An exhaustive search certifies $d=L=12$, so the cap is attained.  


\subsection{Min- and Low-weight Logical Basis}\label{ssec:case-chains}

\paragraph{Greedy Weight Reduction} We can find a min-weight logical basis and its partner logicals via Gaussian elimination followed by a greedy weight reduction, as summarized in Table~\ref{tab:case-bases}.  Both bases attain the code distance, $\wt{\bar X_i}=\wt{\bar Z_i}=12$ for every $i$, but $\mathcal{L}_X^{\min}$ and $\mathcal{L}_Z^{\min}$ are \emph{not} orthonormal partners: their pairing matrix $M_{ij}=\braket{\bar Z_i,\bar X_j}$ has full rank $30$ but is not the identity.  That failure is forced.  Since $H_X=H_Z$, the $X$- and $Z$-logical classes are represented inside the same subspace $\ns H\subset\mathbb{F}_2^{132}$, and every logical has even weight, so $\braket{v,v}=0$ for every representative: a $Z$-logical sharing its partner's support would commute with it, and no minimum-weight basis can be its own symplectic partner.  Orthonormality can be restored by the change of basis $\bar Z'=M^{-1}\bar Z$, which gives $\braket{\bar X_i,\bar Z'_j}=\delta_{ij}$, but only at the cost of weight: the conjugate basis $\mathcal{L}_Z'$ has min/median/max logical weight $18/22/24$.

This greedily generated basis is unstructured: the shift $\sigma$ acts on it as an arbitrary Clifford circuit.  Algorithm~\ref{alg:hypercube-basis} is meant to fix that, but run on the orthonormal pair $(\mathcal{L}_X^{\min},\mathcal{L}_Z')$ that its input hypothesis requires, it emits a massively overcomplete family of $121=11\times11$ logicals spanning a $30$-dimensional space; after the second sector every further orbit of size $11$ adds exactly one new dimension. Since $m=11$ is odd, $R$ is semisimple, and $x^{11}-1=(x-1)g(x)$ with $g(x)=1+x+\dots+x^{10}$.  Here $g$ is irreducible over $\mathbb{F}_2$: the irreducible factors of $x^{11}-1$ are indexed by the cyclotomic cosets of $2$ modulo $11$, and $\ord{2 \bmod 11}=10$ leaves the single nontrivial coset $\{1,2,\dots,10\}$, so $x-1$ and one factor of degree $10$ are all there is.  Hence $R\cong\mathbb{F}_2\times K$ with $K=\mathbb{F}_2[x]/(g)\cong\mathbb{F}_{2^{10}}$.
Reducing $H$ modulo each factor gives $\rank_{\mathbb{F}_2}H(1)=1$, all $55$ rows collapsing onto the all-ones vector of the $12$ blocks, and $\rank_{K}(H\bmod g)=5$; this already accounts for $\rank H=1+5\times10=51$ and hence for $k=132-2\times51=30$.  
\begin{table}[h]
\centering
\begin{tabular}{lccc}
\toprule
basis & $\min$ & median & $\max$ \\
\midrule
$\mathcal{L}_X^{\min}$ & $12$ & $12$ & $12$ \\
$\mathcal{L}_Z^{\min}$ & $12$ & $12$ & $12$ \\
$\mathcal{L}_Z'$ (Orthonormal) & $18$ & $22$ & $24$ \\
\midrule
 $\mathcal{L}^t$, $\mathcal{L}^b$ \ref{fig:case-structure}& $12$ & $12$ & $22$ \\
 ${\mathcal{L}^t}', {\mathcal{L}^b}'$ (Orthonormal) \ref{fig:gatesortho} & $22$ & $24$ & $24$ \\
\bottomrule
\end{tabular}
\caption{Weight spectra of the logical bases of the $[[132,30,12]]$ code, with the
symplectic pairs ordered by weight where a pairing exists. $\mathcal{L}_X^{\min}$ and
$\mathcal{L}_Z'$ form an orthonormal pair, while the structured chains $\mathcal{L}^t$, $\mathcal{L}^b$ and ${\mathcal{L}^t}'$, ${\mathcal{L}^b}'$ has logical shifts and fold-transversal Cliffords (See Figures \ref{fig:case-structure}-\ref{fig:case-structure}).}
\label{tab:case-bases}
\end{table}

\paragraph{Top and Bottom Logical} A easier way to obtain a logical basis is by looking at its top and bottom logicals, which is made especially simple since the top is trivial. First, the top code is simply the all-one matrix, making it dual to the repetition code with exactly weight-2 logicals 
\begin{equation}\label{eq:case-top}
      \mathcal{L}^{t} =(\chi_{i}+\chi_{i+1})\otimes \textbf{1}^{1\times m}\quad, i\in\mathbb{Z}_{12},
\end{equation}
each of weight $2m=22$, where $e_i$ here denote the indicator vector $[\chi_i]_j = \delta _{i,j} $.  Consecutive pairs share a block so they anti-commute, but $ \mathcal{L}^{t}_i$ and $ \mathcal{L}^{t}_{i+2}$ are disjoint split into two sectors of six pairwise disjoint operators $\mathcal{L}^{t}_{e}$ and $ \mathcal{L}^{t}_{o}$, each of rank $5$ since they sum to $\mathbf 1\in\rs H$. Together the two sectors span all shift-invariant logicals. The bottom logical is exactly the latent row, which is also pairwise disjoint and has rank $10$ with the same argument. The remaining bottom logicals come from rank deficiencies of the active rows, which now correspond to the logical operators of the parent matrix, making them logials of the self-dual 2BGA code exactly with abelian quasi-cyclic lifts. 

Let us organize these logicals into five sector disjoint chains, as illustrated in Figure~\ref{fig:case-structure}. Together give $3\times11+2\times6=45$ generators of rank $30$, each with weight $12$ or $22$, and no two of them equal modulo $\rs H$.  Dropping the redundant chain $\mathcal{L}^{b}_{2}$ leaves the $34=11+11+6+6$ generators used in Eq.~\eqref{eq:case-alpha-blocks} and Figure~\ref{fig:case-gates}(c), still of rank $30$.  We note that within a chain the generators have disjoint supports and therefore commute, so the chains are not themselves symplectic bases.

\begin{figure}[t]
\centering
\includegraphics[width=0.95\linewidth]{Assets/case_study_logical.pdf}
\caption{The two logical families of the self-dual $[[132,30,12]]$ GALA code, drawn on the
$12\times11$ (block $\times$ $C_{11}$) grid of its $132$ qubits.  Within one chain the operators
are pairwise disjoint, one color each, and cover every qubit exactly once.
(a) The top logicals.  Each $\mathcal{L}^T_i$ is a pair of complete blocks of weight
$22$; taking every even/odd $i$ makes six of them disjoint, so the twelve modes fall into the two sectors $\mathcal{L}^{t}_{0}$ and $\mathcal{L}^{t}_{1}$, offset by one block. 
(b) The bottom logicals.  The outlined qubits correspond to the representative logical, with orbits under cyclic shift highlighted in different colors.}
\label{fig:case-structure}
\end{figure}

\subsection{Product structure of the chains}\label{ssec:case-products}

Each of the five chains partitions the $132$ qubits, so summing a whole chain gives the
all-ones vector, which is a stabilizer. One consequence is that long products \emph{inside} a chain never need a high-weight representative, since
$\sum_{j\neq j_0}\bar X^{(a)}_j=\mathbf 1+\bar X^{(a)}_{j_0}\equiv\bar X^{(a)}_{j_0}$ modulo
$\rs H$: the product of all but one member of a chain is again weight-$12$.  The redundancy of the third bottom chain extends this \emph{across} chains --- eliminating $\mathcal{L}^{b}_{2}$ against the other two gives
\begin{equation}\label{eq:case-cross}
    \bar X^{(2)}\ =\ (1+x^{8})\,\bar X^{(0)}
    \ +\ (x^{2}+x^{3}+x^{4}+x^{5}+x^{6})\,\bar X^{(1)} ,
\end{equation}
that is, $\bar X^{(2)}_j=\bar X^{(0)}_j+\bar X^{(0)}_{j+8}+\sum_{t=2}^{6}\bar X^{(1)}_{j+t}$ for every $j$, so this particular seven-fold product across $\mathcal{L}^{b}_{0}$ and $\mathcal{L}^{b}_{1}$ also has a weight-$12$ representative.  

\subsection{Transversal Clifford Operations}\label{ssec:case-fold}

With the chains in hand, the dirty cyclic shift can be implemented with a single rigid AOD translation from Section~\ref{sec:logic}; this shifts the bottom chains by one but  fixes the top chain.  The gate is ``dirty'' in the sense that the $11$ sites of a chain carry only $10$
logical qubits, so the $11$-cycle permutes an overcomplete label set rather than a basis.

Secondly, the ZX-duality of this code is $\tau=\mathrm{id}$, so the two fold-transversal operators become to depth-$1$ layers of single-qubit gates:
\begin{equation}\label{eq:case-fold}
    H_\tau=H^{\otimes 132},\qquad S_\tau=(S^{\dagger})^{\otimes 132},
\end{equation}
and $\tau$ is fault tolerant by definition since one qubit gate cannot spread errors. This is special to $\tau=\mathrm{id}$.  For a generic ZX-dual GALA code the rearrangement of Prop.~\ref{def:dual-gala} must actually be performed, and while the resulting $H_\tau$ still only permutes and rotates single qubits, the phase-type gate $S_\tau$ carries a
$CZ$ across every $\tau$-pair, which may drops
fault distance. 


Both logical actions are read off from the gram matrix of the $X$-logical representatives under,
\begin{equation}\label{eq:case-alpha}
    \alpha_{ij}=\braket{\bar X_i,\bar X_j}
    =|\supp{\bar X_i}\cap\supp{\bar X_j}|\ \mathrm{mod}\ 2 .
\end{equation}
Since $\mathbf 1\in\rs H$, every logical has even weight, so $\alpha$ has zero diagonal; and $\alpha$ is nondegenerate, because a $v\in\ns H$ pairing trivially with all of $\ns H$ lies in $(\ns H)^{\perp}=\rs H$ and hence represents the trivial class.  Work in an orthonormal frame, $\braket{\bar X_i,\bar Z_j}=\delta_{ij}$, as supplied by $(\mathcal{L}_X^{\min},\mathcal{L}_Z')$ of Table~\ref{tab:case-bases}, and let $\alpha$ and $\beta$ be the Gram matrices of the $X$- and the $Z$-representatives.  Transversal Hadamard sends $X(v)\mapsto Z(v)$ on the same support, and naming the resulting class by pairing it against the frame gives
\begin{equation}\label{eq:case-H}
    \bar X_i\ \longmapsto\ \prod_j\bar Z_j^{\,\alpha_{ij}},\qquad
    \bar Z_i\ \longmapsto\ \prod_j\bar X_j^{\,\beta_{ij}},\qquad
    \beta_{ij}=\braket{\bar Z_i,\bar Z_j}.
\end{equation}
Dual bases of a nondegenerate symmetric form have mutually inverse Gram matrices, so
$\alpha\beta=I$: transversal $H$ is a logical Hadamard composed with the
$\mathrm{GL}(30,\mathbb{F}_2)$ basis change $\alpha$.  Transversal $S^\dagger$ fixes every
$\bar Z$ and sends
\begin{equation}\label{eq:case-S}
    \bar X(v)\ \longmapsto\ (-1)^{q(v)}\,\bar X(v)\,\bar Z(\alpha v),\qquad
    q(v)=\tfrac12\wt{v}\ \mathrm{mod}\ 2 ,
\end{equation}
the sign being the residual phase $(-i)^{\wt v}=(-1)^{\wt v/2}$ of the calculation above.  So
$S_\tau$ is the diagonal logical Clifford carrying a $\overline{CZ}$ on every $\alpha$-edge and a
single-qubit logical phase wherever $q=1$.  On the weight-$12$ basis of
Table~\ref{tab:case-bases} every $q_i=\tfrac{12}2\bmod 2=0$ and $S_\tau$ is a pure
$\overline{CZ}$ network of $198$ edges; on the chains, the weight-$22$ top modes have
$q=\tfrac{22}2\bmod2=1$ and do carry a phase gate.
\begin{figure}[t]
\centering
\includegraphics[width=0.95\linewidth]{Assets/case_study_gate.pdf}
\caption{Logical gates of the $[[132,30,12]]$ code.
(a) The $\mathbb Z_{11}$ shift: one rigid AOD translation induces
$\bar\sigma:j\mapsto j+1$ simultaneously on the bottom chains while leaving the top chains unchanged.
(b) The logical $\overline{CZ}$ circuit induced by $S_\tau=(S^{\dagger})^{\otimes132}$.  A bundle of four $CZ$s corresponding to the offsets $c(x)=x^2+x^5+x^6+x^9$ repeated at every site. 
On the top chain it is nearest-neighbor coupling corresponding to $12$-cycle
$\Lambda_{12}$.
(c) The anticommutation matrix $\alpha_{ij}=\braket{\bar X_i,\bar X_j}$ over the
$34=6+6+11+11$ generators of the four chains $\mathcal{L}^{t}_{0}$, $\mathcal{L}^{t}_{1}$,
$\mathcal{L}^{b}_{0}$, $\mathcal{L}^{b}_{1}$, that is, over the $45$ generators of
Section~\ref{ssec:case-chains} with the redundant third bottom chain $\mathcal{L}^{b}_{2}$
dropped; the remaining $34$ still have rank $30$ modulo $\rs H$.}
\label{fig:case-gates}
\end{figure}

The chains make the structure of $\alpha$ explicit (Figure~\ref{fig:case-gates}(b,c)), and in
the labeling of Section~\ref{ssec:case-chains} the two families come out with the same
shape.  Since $\sigma$ is a permutation it preserves the pairing, and the chain generators are
$\sigma$-orbits, so $\alpha_{(a,j),(b,j')}$ depends only on $j'-j$: every block of $\alpha$ is
circulant.  In the ordering
$(\mathcal{L}^{t}_{0},\mathcal{L}^{t}_{1},\mathcal{L}^{b}_{0},\mathcal{L}^{b}_{1})$
\begin{equation}\label{eq:case-alpha-blocks}
    \alpha\ =\
    \begin{pmatrix}
        0 & 1+y^{-1}\\ 1+y & 0
    \end{pmatrix}
    \ \oplus\
    \begin{pmatrix}
        0 & c(x)\\ c(x) & 0
    \end{pmatrix},
    \qquad
    c(x)=x^{2}+x^{5}+x^{6}+x^{9},
\end{equation}
where $y$ shifts the block index inside a top sector and $c$ is symmetric.  The top block is the adjacency matrix $\Lambda_{12}$ of the
$12$-site cycle, written in the two-sector ordering: $\bar D_b$ and $\bar D_{b'}$ share an
odd number of qubits only when $b'=b\pm1$, and the cycle alternates between
$\mathcal{L}^{t}_{0}$ and $\mathcal{L}^{t}_{1}$, so it is bipartite with the two offsets
$\{0,-1\}$.  


Choosing an orthonormal basis instead of the chains turns the same gate into a perfect
matching: $\alpha$ decomposes into $15$ disjoint pairs, which is the $30$ logical
qubits, split as $10$ pairs across the two bottom chains and $5$ pairs across the two top sectors.  In that basis $S_\tau$ is $15$ disjoint $\overline{CZ}$ gates together with phases on $13$ of the $30$ logical qubits, and $H_\tau$ is $15$ copies of $(\bar H\otimes\bar H)\cdot\overline{\mathrm{SWAP}}$.  That basis is the legible one, but its representatives run from weight $12$ up to $24$; the chains keep every representative at $12$ or $22$, so they are the ones to compile against. See Figure~\ref{fig:gatesortho}.


\begin{figure}
    \centering
    \includegraphics[width=0.95\linewidth]{Assets/case_study_orthonormal.pdf}
    \caption{The logical basis with higher weight but simple transversal Cliffords. (a) the logical basis--the bottom logicals have higher weight compared to baseline, but the logical shifts have a clean description. (b) the transversal Hadamard implements a logical Hadamard gate followed by a logical SWAP between disjoint sectors. (c) transversal phase gate implements a CZ across disjoint sectors.}
    \label{fig:gatesortho}
\end{figure}